\documentclass[10pt,a4paper]{amsart}
\usepackage[margin=19mm]{geometry}
\usepackage{amsmath,amssymb,mathtools}
\usepackage[scr=rsfs,cal=euler]{mathalfa}
\usepackage{xfrac}
\usepackage[unicode,hidelinks,bookmarks=false]{hyperref}
\newcommand{\ie}{i.e.\ }
\newcommand{\cf}{cf.\ }
\newcommand{\resp}{respectively\ }
\newcommand{\Subsection}{\subsection}
\providecommand{\nobreakdash}{\nobreak\hbox{-}}
\setlength{\emergencystretch}{2em}
\multlinegap0pt
\usepackage{amssymb}
\newtheorem{proposition}{Proposition}
\title{On the differential equations of frozen Calogero-Moser-Sutherland particle models}
\author{Michael Voit}
\date{Source: arXiv:2312.02685v1 (CC BY 4.0)}
\begin{document}
\maketitle

\noindent\emph{Source and licence.} On the differential equations of frozen Calogero-Moser-Sutherland particle models. Version arXiv:2312.02685v1 (CC BY 4.0), distributed by arXiv under the Creative Commons Attribution 4.0 International licence, http://creativecommons.org/licenses/by/4.0/, as recorded in the arXiv licence metadata for this exact version and shown on the abstract page https://arxiv.org/abs/2312.02685v1. Full licence notice: \texttt{licenses/arxiv-2312.02685-CC-BY-4.0.txt}.\par\smallskip

\noindent\textit{Editorial scope.} Counted unit: the article's proposition decreasing-error-jacobi-stat-trigono (source0.tex lines 1150--1155). The article's complete original proof of it is delivered with it (source0.tex lines 1157--1220, one \texttt{proof} environment of 64 lines). No mathematics is written, repaired or re-derived: the statement and the proof are the article's own lines, in the article's own notation, and no retained line is substituted or re-flowed.  Retained uncounted as the source-local context the proof needs: lines 292--295; lines 353--357; lines 762--768; lines 788--792; lines 1069--1074; lines 1132--1139.  Nothing was omitted from any retained span, so the package carries no omission record.  No retained line contains a citation, so the wrapper carries no bibliography and no bibliography entry.  No retained line contains a figure environment, an image-inclusion command or drawing code, so no image is delivered and the package declares no asset dependency.  Editorial formatting only, in addition: an AMS \texttt{amsart} preamble that restates the article's own declarations and macros used by the retained lines, namely the environments \texttt{proposition} on separate counters, together with the standard packages those lines need.  The retained environments print in this document as Proposition 1 in order of appearance, whereas the article prints the same items with the numbering of its own full document.  The complete native source of the article as distributed in the arXiv e-print for this exact version is bundled unchanged as \texttt{sources/151032/source0.tex}.
\begin{equation}\label{ODE-tilde-a}
{\psi}_i^\prime(t) = \sum_{j: j\ne i}  \frac{1}{\psi_{i}(t)-\psi_{j}(t)}
-\frac{ N(N-1)}{2} \psi_{i}(t)\quad(i=1,\ldots,N)
\end{equation}

\begin{proposition}\label{decreasing-error-a-stat}
  Let  $\psi(t),\tilde \psi(t)$ be solutions of (\ref{ODE-tilde-a}) with $\psi(0),\tilde \psi(0)\in C_N^A$.
Then for $t\ge0$,
  $$\Bigl\|\psi(t)-\tilde \psi(t)\Bigr\|\le  e^{- \frac{N(N-1)}{2} t}\Bigl\|\psi(0)-\tilde \psi(0)\Bigr\|.$$ 
 \end{proposition}

\begin{align}\label{ODE-tilde-b}
{\psi}_i^\prime(t) = &\sum_{j\ne i}  \frac{1}{\psi_{i}(t)-\psi_{j}(t)}+
\sum_{j\ne i}  \frac{1}{\psi_{i}(t)+\psi_{j}(t)}
 \notag\\
&\quad
+\frac{ \nu}{  \psi_{i}(t)} -   N(N+\nu-1)\cdot \psi_{i}(t)
\end{align}

\begin{proposition}\label{decreasing-error-b-stat}
  Let  $\psi(t),\tilde \psi(t)$ be solutions of (\ref{ODE-tilde-b}) with $\psi(0),\tilde \psi(0)\in C_N^A$.
Then for $t\ge0$,
  $$\Bigl\|\psi(t)-\tilde \psi(t)\Bigr\|\le  e^{- N(N+\nu-1) t}\Bigl\|\psi(0)-\tilde \psi(0)\Bigr\|.$$ 
 \end{proposition}

\begin{equation}\label{ODE-jacobi}
\frac{d}{dt}x_i(t)
		=(p-q)-(p+q)x_i(t)
			+2\sum_{j:\> j\neq i}\frac{1-x_i(t)
			x_j(t)}{x_i(t)-x_j(t)}\quad ( i=1,\dots,N)
\end{equation}

   \begin{equation}\label{trig-ode}
        \begin{split}
        \frac{d}{dt} \tau_i(t)=&
           (q-p)\cot\left(\frac{\tau_i(t)}{2}\right)+2(p+1-N) \cot(\tau_i(t))\\
          &+\sum_{j: j\ne i}\left(\cot\left(\frac{\tau_i(t)-\tau_j(t)}{2}\right)+
          \cot\left(\frac{\tau_i(t)+\tau_j(t)}{2}\right)\right).
          \end{split}
          \end{equation}

 \begin{proposition}\label{decreasing-error-jacobi-stat-trigono}
  Let  $\tau(t),\tilde \tau(t)$ be solutions of (\ref{trig-ode}) with $\tau(0),\tilde \tau(0)\in \tilde A_N$.
Then for $t\ge0$,
$$\Bigl\|\tau(t)-\tilde \tau(t)\Bigr\|\le  e^{- c  t}\Bigl\|\tau(0)-\tilde \tau(0)\Bigr\|  \quad\text{with}\quad
c=\frac{p+q+2\cdot min(p,q)+2-2N}{4}>\frac{N-1}{2}\ge0.$$ 
  \end{proposition}

 \begin{proof} We proceed as for  Propositions \ref{decreasing-error-a-stat} and \ref{decreasing-error-b-stat}.
   Let $r_j(t):=\tau_j(t)-\tilde \tau_j(t)$ for $j=1,\ldots,N$, and $D(t):=\frac{1}{2}\sum_{j=1}^N r_j(t)^2$.
   Then,  by (\ref{trig-ode}),
\begin{equation}
  D^\prime(t) = 
  \sum_{j=1}^N\Bigl( \tau_j^\prime(t)- \tilde \tau_j^\prime(t)\Bigr)\Bigl( \tau_j(t)- \tilde \tau_j(t)\Bigr)
=:  A_1(t)+A_2(t) \end{equation}
with
\begin{align}
A_1(t)= \sum_{l,j: \> l\ne j}\Bigl( \tau_j(t)- \tilde \tau_j(t)\Bigr)\Bigl(&\cot\Bigl(\frac{\tau_j(t)-\tau_l(t)}{2}\Bigr)-
\cot\Bigl(\frac{\tilde \tau_j(t)-\tilde \tau_l(t)}{2}\Bigr)\Bigr)\notag\\
&+\Bigl(\cot\Bigl(\frac{\tau_j(t)+\tau_l(t)}{2}\Bigr)-
\cot\Bigl(\frac{\tilde \tau_j(t)+\tilde \tau_l(t)}{2}\Bigr)\Bigr)\notag
\end{align}
and
\begin{align}
  A_2(t)=\sum_{i=1}^N \Bigl( \tau_j(t)- \tilde \tau_j(t)\Bigr)
  \cdot\Bigl( &(q-p)\Bigl(\cot\left(\frac{\tau_i(t)}{2}\right)-\cot\left(\frac{\tilde\tau_i(t)}{2}\right)\Bigr)\notag\\
&+2(p+1-N)(\cot\tau_i(t) -\cot\tilde\tau_i(t))\Bigr).
\notag
\end{align}
In order to estimate  $A_2(t)$ we conclude from the mean value theorem that
\begin{align}\label{sine-estimate}
(q-p)&\Bigl(\cot\Bigl(\frac{\tau_i(t)}{2}\Bigr)-\cot\Bigl(\frac{\tilde\tau_i(t)}{2}\Bigr)\Bigr)
+2(p+1-N)(\cot\tau_i(t) -\cot\tilde\tau_i(t))\notag\\
&= -(\tau_i(t)-\tilde\tau_i(t))\Bigl(\frac{q-p}{2\sin^2 (\hat\tau_i(t)/2)} + \frac{2(p+1-N)}{\sin^2 (\hat\tau_i(t))}\Bigr)
\end{align}
with some $\hat\tau_i(t)$ between $\tau_i(t)$ and  $\tilde\tau_i(t)$. This implies for $q\ge p$ that
\begin{equation}\label{est-jacobi-a2}  A_2(t)\le -\frac{q+3p+4-4N}{2}\sum_{i=1}^N \Bigl( \tau_j(t)- \tilde \tau_j(t)\Bigr)^2.
\end{equation}
In order to handle  $A_1(t)$, we write it as
\begin{align}
 & A_1(t)= \notag\\ =& 
  \sum_{l,j: \> l> j}\Bigl( \tau_j(t)+ \tau_l(t) -\tilde \tau_j(t)-\tilde \tau_l(t)\Bigr)
 \Bigl(\cot\Bigl(\frac{\tau_j(t)+\tau_l(t)}{2}\Bigr)-
\cot\Bigl(\frac{\tilde \tau_j(t)+\tilde \tau_l(t)}{2}\Bigr)\Bigr)\notag\\
+&
  \sum_{l,j: \> l> j}\Bigl( \tau_j(t)- \tau_l(t) -\tilde \tau_j(t)+\tilde \tau_l(t)\Bigr)
 \Bigl(\cot\Bigl(\frac{\tau_j(t)-\tau_l(t)}{2}\Bigr)-
\cot\Bigl(\frac{\tilde \tau_j(t)-\tilde \tau_l(t)}{2}\Bigr)\Bigr).\notag
\end{align}
Hence, again by the mean value theorem for the cotangens, for some $\tau_{j,l,1},\tau_{j,l,2}\in [0,\pi]$,
\begin{align}
  A_1(t)= & -\frac{1}{2} \sum_{l,j: \> l>j}\Bigl( \tau_j(t)+ \tau_l(t) -\tilde \tau_j(t)-\tilde \tau_l(t)\Bigr)^2\frac{1}{\sin^2( \tau_{j,l,1})}
\notag\\
&-\frac{1}{2} \sum_{l,j: \> l>j}\Bigl( \tau_j(t)- \tau_l(t) -\tilde \tau_j(t)+\tilde \tau_l(t)\Bigr)^2\frac{1}{\sin^2( \tau_{j,l,2})}\notag\\
\le  -\frac{1}{4} &\sum_{l,j: \> l\ne j}\Bigl(\Bigl( \tau_j(t)+ \tau_l(t) -\tilde \tau_j(t)-\tilde \tau_l(t)\Bigr)^2
+\Bigl( \tau_j(t)- \tau_l(t) -\tilde \tau_j(t)+\tilde \tau_l(t)\Bigr)^2\Bigr).\notag
\end{align}
As
\begin{align}
  \Bigl( \tau_j(t)+ &\tau_l(t) -\tilde \tau_j(t)-\tilde \tau_l(t)\Bigr)^2+
  \Bigl( \tau_j(t)-\tau_l(t) -\tilde \tau_j(t)+\tilde \tau_l(t)\Bigr)^2\notag\\
&= 2(\tau_j(t)-\tilde \tau_j(t))^2 +2(\tau_l(t)-\tilde \tau_l(t))^2,\end{align}
we conclude that
\begin{equation}\label{est-jacobi-a1}
  A_1(t)\le -(N-1)\sum_{j=1}^N (\tau_j(t)-\tilde \tau_j(t))^2.\end{equation}
In summary we obtain from (\ref{est-jacobi-a1}) and (\ref{est-jacobi-a2}) for $q\ge p$ that
$$D^\prime(t) \le   -\frac{q+3p+2-2N}{2} D(t).$$
This and the lemma of Gronwall yield the claim for $q\ge p$.

For the case $q\le p$ we notice that the ODE (\ref{ODE-jacobi}) is invariant under the transform $x(t)\mapsto -x(t)$
when the parameters $p$ and $q$ are interchanged. This leads to the theorem for $q\le p$.
 \end{proof}

\end{document}
